\chapter{Auxiliary processes and martingales}
\label{chap:auxiliary}

The analysis of the $\aRTS$ family rests on a handful of auxiliary stochastic
processes built from the assignments and responses. This chapter collects their
definitions and elementary (mostly martingale) properties. It corresponds to the
setup of the proof sketch (Section~3.2) and of Appendix~A.

\section{Conditional probabilities and the assignment martingale}
\label{sec:aux_assignment}

\begin{definition}[Conditional selection probability]\label{def:cond_prob}
\uses{def:assignment, def:filtration}
For $k \in [K]$ and $m \ge 1$,
\[
  p_{m,k} := \E[X_{m,k} \mid \filt_{m-1}].
\]
For a RAR procedure these coincide with the allocation probabilities: with the
indexing convention of Definition~\ref{def:rar}, $p_{m,k}$ is the $k$-th
component of the vector used to draw $X_m$.
\end{definition}

\begin{definition}[Assignment martingale]\label{def:M}
\lean{AlphaRAR.assignMG}\leanok
\uses{def:cond_prob}
For $k \in [K]$ set $\Delta M_{m,k} := X_{m,k} - p_{m,k}$ and
\[
  M_{n,k} := \sum_{m=1}^n (X_{m,k} - p_{m,k}), \qquad M_n := (M_{n,1},\dots,M_{n,K}).
\]
\end{definition}

\begin{lemma}[$M$ is a martingale]\label{lem:M_martingale}
\lean{AlphaRAR.martingale_assignMart}\leanok
\uses{def:M, def:filtration}
For each $k \in [K]$, $(M_{n,k})_{n\ge0}$ is an $(\filt_n)$-martingale with
$M_{0,k}=0$ and bounded increments $|\Delta M_{m,k}| \le 1$.
\end{lemma}

\begin{proof}\leanok
\uses{def:cond_prob}
$\E[\Delta M_{m,k}\mid\filt_{m-1}] = \E[X_{m,k}\mid\filt_{m-1}] - p_{m,k} = 0$
by Definition~\ref{def:cond_prob}, and $X_{m,k},p_{m,k}\in[0,1]$.
\end{proof}

\begin{lemma}[The assignment martingale is $O_p(\sqrt n)$]\label{lem:M_Op}
\lean{AlphaRAR.isBigOpOne_assignMart_div_sqrt}\leanok
\uses{def:M, lem:M_martingale}
For each $k \in [K]$, $M_{n,k} = O_p(\sqrt n)$.
\end{lemma}

\begin{proof}\leanok
\uses{lem:mart_bdd_incr_Op}
$M_{\cdot,k}$ is a martingale with $M_{0,k}=0$ and increments bounded by $1$
(Lemma~\ref{lem:M_martingale}), so Lemma~\ref{lem:mart_bdd_incr_Op} applies with
$c = 1$.
\end{proof}

\begin{lemma}[Count decomposition]\label{lem:count_decomp}
\lean{AlphaRAR.count_eq}\leanok
\uses{def:M, def:counts}
For all $n\ge1$ and $k\in[K]$,
$\displaystyle N_{n,k} = \sum_{m=1}^n p_{m,k} + M_{n,k}$.
\end{lemma}

\begin{proof}\leanok
Sum the identity $X_{m,k} = p_{m,k} + \Delta M_{m,k}$ over $m \le n$.
\end{proof}

\section{The response martingale}
\label{sec:aux_response}

\begin{definition}[Centered response martingale]\label{def:Q}
\lean{AlphaRAR.respMG, AlphaRAR.respMart}\leanok
\uses{def:assignment, def:responses}
For $k \in [K]$ set $Y_{m,k} := \xi_{m,k} - \theta_k$ and
\[
  Q_{n,k} := \sum_{m=1}^n X_{m,k}\,(\xi_{m,k} - \theta_k)
           = \sum_{m=1}^n X_{m,k}\,Y_{m,k},
  \qquad Q_n := (Q_{n,1}, \dots, Q_{n,K}).
\]
\end{definition}

\begin{lemma}[$Q$ is a martingale]\label{lem:Q_martingale}
\lean{AlphaRAR.martingale_respMart}\leanok
\uses{def:Q, def:filtration}
For each $k \in [K]$, if the responses are integrable then $(Q_{n,k})_{n\ge0}$ is an
$(\filt_n)$-martingale.
\end{lemma}

\begin{proof}\leanok
Write $\Delta Q_{n,k} = X_{n,k}(\xi_{n,k} - \theta_k)$, where $X_{n,k}$ is the
indicator that the assignment $A_n$ equals arm $k$. The subtle point is the
\emph{filtration convention}: $X_{n,k}$ is measurable with respect to the
action-augmented $\sigma$-algebra $\mathcal{G}_n := \filt_{n-1}\vee\sigma(A_n)$
(\texttt{filtrationAction}) but \emph{not} with respect to $\filt_{n-1}$ --- indeed,
if it were, the assignment martingale $M_{n,k}$ of Definition~\ref{def:M} would be
identically zero. Conditioning on $\mathcal{G}_n$, the response $\xi_{n,k}$ is drawn
from the arm's reward kernel independently of the past given the arm, so its
conditional mean is $\theta_k$; pulling out the $\mathcal{G}_n$-measurable factor
$X_{n,k}$ gives $\E[\Delta Q_{n,k}\mid\mathcal{G}_n] = X_{n,k}(\theta_k-\theta_k)=0$.
The tower property ($\filt_{n-1}\subseteq\mathcal{G}_n$) yields
$\E[\Delta Q_{n,k}\mid\filt_{n-1}] = 0$.

This is formalized in the \texttt{LeanMachineLearning} algorithm--environment
framework: the arm is the action and the response is the feedback of a stationary
environment with per-arm reward kernel $\nu$, and $\theta_k = \int x\,\nu_k(dx)$ is
the kernel mean. The conditional-mean computation is \texttt{condExp\_feedback}, and
the two filtrations $\filt$ and $\mathcal{G}$ are \texttt{filtration} and
\texttt{filtrationAction}.
\end{proof}

\begin{lemma}[Quadratic variation of $Q$]\label{lem:Q_quad_var}
\uses{def:Q, def:counts}
Under Condition~\textbf{A}, the predictable quadratic variation of $Q_{\cdot,k}$
is $\langle Q_{\cdot,k}\rangle_n = V_k N_{n,k}$, and for $k \ne j$ the cross
variation $\langle Q_{\cdot,k}, Q_{\cdot,j}\rangle_n = 0$.
\end{lemma}

\begin{proof}
\uses{lem:Q_martingale, lem:qv_incr, lem:qv_orthogonal}
Sum the conditional second moments of Lemma~\ref{lem:Q_martingale} using
$\sum_{m\le n} X_{m,k} = N_{n,k}$. The cross terms vanish because
$X_{m,k}X_{m,j} = 0$ for $k \ne j$ (each patient is assigned to one arm).
\end{proof}

\begin{lemma}[Estimator error via $Q$]\label{lem:theta_error_Q}
\lean{AlphaRAR.theta_error_Q}\leanok
\uses{def:estimator, def:Q, def:bahadur}
On $\{N_{m,k}\ge1\}$, the leading term of the estimator error is
$\dfrac{1}{N_{m,k}}\sum_{j=1}^m X_{j,k}\xi_{j,k} - \theta_k
 = \dfrac{Q_{m,k}}{N_{m,k}}$.
Consequently, under the Bahadur representation
(Definition~\ref{def:bahadur}),
$\hat{\theta}_{m,k} - \theta_k = \dfrac{Q_{m,k}}{N_{m,k}} + o(N_{m,k}^{-1/2})$.
\end{lemma}

\begin{proof}
$\frac{1}{N_{m,k}}\sum_j X_{j,k}\xi_{j,k} - \theta_k
= \frac{1}{N_{m,k}}\sum_j X_{j,k}(\xi_{j,k}-\theta_k)
= \frac{Q_{m,k}}{N_{m,k}}$, using $\sum_j X_{j,k} = N_{m,k}$.
\leanok
\end{proof}

\section{The auxiliary process \texorpdfstring{$U$}{U} and the last hitting time}
\label{sec:aux_U}

\begin{definition}[Auxiliary process $U$]\label{def:U}
\lean{AlphaRAR.auxU}\leanok
\uses{def:M, def:plugin_target}
For a fixed $\alpha \in [0,1)$, $k \in [K]$ and $n \ge 1$,
\[
  U_{n,k} := \sum_{m=1}^{n-1} \alpha\, \hat{\rho}_{m,k} + M_{n,k} - n\, \hat{\rho}_{n,k}.
\]
\end{definition}

\begin{lemma}[Increment of $U$]\label{lem:U_increment}
\lean{AlphaRAR.auxU_succ_sub}\leanok
\uses{def:U}
For all $n \ge 1$ and $k \in [K]$,
\[
  \Delta U_{n,k} = U_{n,k} - U_{n-1,k}
  = \alpha\, \hat{\rho}_{n-1,k} - p_{n,k} + \Delta\big(N_{n,k} - n\hat{\rho}_{n,k}\big),
\]
where we used $\Delta M_{n,k} = X_{n,k} - p_{n,k}$ and $\Delta N_{n,k} = X_{n,k}$.
\end{lemma}

\begin{proof}\leanok
Direct computation from Definition~\ref{def:U}, expanding
$M_{n,k} - M_{n-1,k} = X_{n,k} - p_{n,k}$ and telescoping.
\end{proof}

\begin{definition}[Last under-sampling time]\label{def:hitting}
\lean{AlphaRAR.hitting}\leanok
\uses{def:counts, def:plugin_target}
For $k \in [K]$ and $n \ge 1$, the \emph{last under-sampling time} before $n$ is
\[
  \ell_{n,k} := \max\Big\{\, m \le n : \frac{N_{m,k}}{m} \le \hat{\rho}_{m,k} \,\Big\}.
\]
It is the last time up to $n$ at which arm $k$ is under-sampled (its proportion
is at most its estimated target).
\end{definition}

\begin{lemma}[Basic properties of the hitting time]\label{lem:hitting_basic}
\lean{AlphaRAR.hitting_basic}\leanok
\uses{def:hitting}
For each $k$, $(\ell_{n,k})_{n\ge1}$ is non-decreasing and $\ell_{n,k} \le n$.
\end{lemma}

\begin{proof}\leanok
Immediate from the definition: enlarging $n$ enlarges the set over which the max
is taken, and the max is bounded by $n$.
\end{proof}

\textbf{Lean remark.} Well-definedness of $\ell_{n,k}$ requires the set
$\{m\le n : N_{m,k}/m \le \hat\rho_{m,k}\}$ to be nonempty; a convenient
convention (used implicitly in the paper) is that it contains a base time such
as the end of the burn-in, or one sets $\ell_{n,k}=0$ when the set is empty.
The proofs only use Lemma~\ref{lem:hitting_basic} together with the sign
property of the next lemma.

\begin{lemma}[Sign at the hitting time]\label{lem:hitting_sign}
\lean{AlphaRAR.hitting_sign}\leanok
\uses{def:hitting, def:aRTS}
For $m > \ell_{n,k}$ (and $m \le n$), $\frac{N_{m,k}}{m} > \hat{\rho}_{m,k}$;
in particular, for an $\aRTS$ design with $m \ge Km_0$, $p_{m+1,k} \le \alpha\hat{\rho}_{m,k}$.
\end{lemma}

\begin{proof}\leanok
By maximality of $\ell_{n,k}$, no $m$ with $\ell_{n,k} < m \le n$ satisfies
$N_{m,k}/m \le \hat{\rho}_{m,k}$. The second statement is the throttling
condition \eqref{eq:throttle}.
\end{proof}
